\section{Conclusiones}
\label{sec:conclusiones}

\begin{table}[htbp]
    \centering
    \caption{Conclusiones del an\'{a}lisis de flexibilidad de la l\'{i}nea SIM-002.}
    \label{tab:conclusiones}
    \setcounter{itemcount}{0}
    \begin{tabularx}{\textwidth}{@{}NX@{}}
        \toprule
        \multicolumn{1}{c}{\textbf{Item}} & \textbf{Conclusion} \\
        \midrule
        & La tuber\'{i}a de descarga de la bomba de recirculaci\'{o}n 1 del sistema DDW (SIM-002) cumple los criterios de esfuerzo de ASME B31.3-2016 en los nueve casos de carga evaluados: la evaluaci\'{o}n de cumplimiento de c\'{o}digo del modelo arroj\'{o} CODE COMPLIANCE EVALUATION PASSED. \\
        & El esfuerzo de c\'{o}digo m\'{a}ximo es \num{39668.2}~\si{kPa} en el nodo 78, caso 6 (SUS) W+P2, con una relaci\'{o}n de utilizaci\'{o}n de 34.5~\% frente al admisible de \num{115142.4}~\si{kPa}; la condici\'{o}n que gobierna el dise\~{n}o es la sostenida (peso m\'{a}s presi\'{o}n), no el rango de expansi\'{o}n t\'{e}rmica. \\
        & Los desplazamientos m\'{a}ximos del sistema son 5.0~\si{mm} en las componentes horizontales (cierre de la holgura de los soportes \textit{w/gap} de los nodos 10 y 170) y 2.0~\si{mm} en la vertical, y las cargas en soportes son moderadas (carga vertical m\'{a}xima \num{5529}~\si{N} en el soporte del nodo 70 en operaci\'{o}n), por lo que la configuraci\'{o}n de soporter\'{i}a del isom\'{e}trico es adecuada. \\
        & La comparaci\'{o}n de las cargas en la boquilla de la bomba con los admisibles del fabricante del equipo queda pendiente como verificaci\'{o}n externa, por no disponerse de dichos admisibles a la fecha. \\
        \bottomrule
    \end{tabularx}
\end{table}

En s\'{i}ntesis, la l\'{i}nea SIM-002 es t\'{e}cnicamente viable seg\'{u}n
ASME B31.3-2016 con una utilizaci\'{o}n m\'{a}xima de 34.5~\% del admisible. No se identifican riesgos de sobre-esfuerzo, interferencias por
expansi\'{o}n ni sobrecarga de soportes, por lo que no se requieren acciones
correctivas sobre la configuraci\'{o}n analizada.
